% =====================================================================
\chapter{Soluciones}
\label{ap:soluciones}
% =====================================================================

Las referencias a los apuntes siguen el formato «M05, §3.2»: módulo, y número de tema y sección tal como aparecen en el PDF de cada módulo. En los pasos intermedios se conservan al menos 4 cifras significativas y los resultados se dan con 3, como piden las instrucciones; cuando existe, se da también el valor exacto como fracción.

% ---------------------------------------------------------------------
\section{Parte A: clave de respuestas}

\begin{center}
\small
\renewcommand{\arraystretch}{1.12}
\begin{tabularx}{\textwidth}{@{}c c >{\raggedright\arraybackslash}X >{\raggedright\arraybackslash}p{2.6cm}@{}}
\toprule
Pregunta & Resp. & Justificación & Apuntes \\
\midrule
1 & (b) & Se obtiene aplicando repetidamente $P(A\cap B)=P(A)\,P(B\mid A)$, sin ninguna aproximación. & M01, §2.2 \\
2 & (c) & Dos tokens de contexto y uno predicho: $28^{3}=21\,952$ casillas. & M01, §2.3 \\
3 & (c) & $\PPL=\exp(-\ln\tfrac14)=4$: el modelo duda por igual entre 4 opciones. & M01, §5.1 \\
4 & (b) & Cada $z_i/T\to 0$, luego cada $e^{z_i/T}\to 1$ y el cociente tiende a $1/K$. & M01, §4.1 \\
5 & (c) & $\cos=3/(\sqrt2\cdot 3)=1/\sqrt2\approx 0{,}707$. & M02, §2.3 \\
6 & (b) & El único~1 de $e_i$ selecciona la fila $i$: es la capa de embeddings. & M02, §5.1 \\
7 & (a) & \texttt{@} actúa sobre las dos últimas dimensiones: $(128,512)(512,50\,000)$ para cada uno de los 8 elementos del lote. & M03, §6.2 \\
8 & (b) & La segunda fila es el doble de la primera: $\det=0$, rango~1 y $\lambda=0$ es autovalor (el otro es la traza, 4). & M03, §1.3 y §3.1 \\
9 & (c) & Eckart--Young: $\big(\sum_{i>1}\sigma_i^{2}\big)^{1/2}=\sqrt{9+1}=\sqrt{10}$. & M03, §5.1 \\
10 & (b) & Sin centrar, $X\T X/n$ es la matriz de segundos momentos y su autovector dominante se orienta hacia la media. & M03, §5.4 \\
11 & (c) & Las formas se alinean por la derecha: $128$ frente a $512$ no coinciden y ninguna es~1. & M03, §6.3 \\
12 & (b) & $\partial_x f=2xy=6$ y $\partial_y f=x^{2}=1$ en $(1,3)$. & M04, §1.2 \\
13 & (c) & Teorema del gradiente de la entropía cruzada; se deriva en el problema B.4(a). & M04, §3.4 \\
14 & (b) & En el desarrollo de Taylor se cancelan los términos pares y queda $\tfrac{h^{2}}{6}f'''(\xi)$. & M04, §6.1 \\
15 & (b) & $\Var(X)=p(1-p)=0{,}2\cdot 0{,}8=0{,}16$. & M05, §2.3 \\
16 & (c) & Maximizar $\sum_j N_j\ln p_j$ con $\sum_j p_j=1$ da contar y normalizar; (b) es el suavizado de Laplace. & M05, §4.2; M01, §3.2 \\
17 & (b) & La desviación típica de la media de $B$ ejemplos es $\sigma/\sqrt B$, y $\sqrt4=2$. & M05, §5.6 \\
18 & (c) & $H=\log_2 8=3$ bits. & M06, §1.2 \\
19 & (c) & Desigualdad de Gibbs; la divergencia no es simétrica. & M06, §2.2 \\
20 & (a) & $H(p,q)=H(p)+\KL{p}{q}$, luego $\KL{p}{q}=2{,}5-2=0{,}5$ bits. & M06, §2.3 \\
\bottomrule
\end{tabularx}
\end{center}

% ---------------------------------------------------------------------
\section{Parte B: soluciones de los problemas}

% =====================================================================  B.1
\begin{solucion}{eje:p1}
\apartado{(a)}
\datos corpus «mar», «ama», «rama»; $\V=\{\fin,\texttt{a},\texttt{m},\texttt{r}\}$.
\planteamiento añadir \fin\ al principio y al final de cada palabra, contar cada par (anterior, siguiente) y dividir cada fila por su total (M01, §3.2).
\pasos
\begin{enumerate}[label=\arabic*.]
  \item Secuencias con marcadores: \fin mar\fin, \fin ama\fin, \fin rama\fin.
  \item Bigramas de \fin mar\fin: $(\fin,\texttt m)$, $(\texttt m,\texttt a)$, $(\texttt a,\texttt r)$, $(\texttt r,\fin)$; son~4.
  \item Bigramas de \fin ama\fin: $(\fin,\texttt a)$, $(\texttt a,\texttt m)$, $(\texttt m,\texttt a)$, $(\texttt a,\fin)$; son~4.
  \item Bigramas de \fin rama\fin: $(\fin,\texttt r)$, $(\texttt r,\texttt a)$, $(\texttt a,\texttt m)$, $(\texttt m,\texttt a)$, $(\texttt a,\fin)$; son~5.
  \item Total: $4+4+5=13$. Conteos y, entre paréntesis, probabilidades (fila: anterior; columna: siguiente):
\end{enumerate}
\begin{center}
\small
\begin{tabular}{@{}lccccc@{}}
\toprule
 & \fin & \texttt a & \texttt m & \texttt r & Total \\
\midrule
\fin & $0$ & $1\ (1/3)$ & $1\ (1/3)$ & $1\ (1/3)$ & $3$ \\
\texttt a & $2\ (0{,}4)$ & $0$ & $2\ (0{,}4)$ & $1\ (0{,}2)$ & $5$ \\
\texttt m & $0$ & $3\ (1)$ & $0$ & $0$ & $3$ \\
\texttt r & $1\ (0{,}5)$ & $1\ (0{,}5)$ & $0$ & $0$ & $2$ \\
\bottomrule
\end{tabular}
\end{center}
\resultado \textbf{13 bigramas; $N$ y $P$ como en la tabla.}
\comprobacion los totales de fila suman $3+5+3+2=13$ y cada fila de $P$ suma~1.
\criterio 2 puntos la matriz $N$; 1 punto la normalización por filas.

\apartado{(b)}
\datos la matriz $P$ del apartado (a).
\planteamiento regla de la cadena con la hipótesis de Markov de orden~1: la probabilidad de una palabra es el producto de $P(\text{siguiente}\mid\text{anterior})$ a lo largo de \fin palabra\fin\ (M01, §2.3).
\pasos
\begin{enumerate}[label=\arabic*.]
  \item $P(\text{mar})=P(\texttt m\mid\fin)P(\texttt a\mid\texttt m)P(\texttt r\mid\texttt a)P(\fin\mid\texttt r)=\tfrac13\cdot1\cdot\tfrac15\cdot\tfrac12=\tfrac1{30}\approx 0{,}03333$.
  \item $P(\text{rama})=\tfrac13\cdot\tfrac12\cdot\tfrac25\cdot1\cdot\tfrac25=\tfrac{2}{75}\approx 0{,}02667$.
  \item $P(\text{mama})=P(\texttt m\mid\fin)P(\texttt a\mid\texttt m)P(\texttt m\mid\texttt a)P(\texttt a\mid\texttt m)P(\fin\mid\texttt a)=\tfrac13\cdot1\cdot\tfrac25\cdot1\cdot\tfrac25=\tfrac{4}{75}\approx 0{,}05333$.
  \item $P(\text{ama})=P(\texttt a\mid\fin)P(\texttt m\mid\texttt a)P(\texttt a\mid\texttt m)P(\fin\mid\texttt a)=\tfrac13\cdot\tfrac25\cdot1\cdot\tfrac25=\tfrac{4}{75}$.
  \item «mama» solo usa transiciones observadas. Como el modelo solo mira el token anterior y $P(\texttt a\mid\texttt m)=1$, el prefijo «ma» cuesta $P(\texttt m\mid\fin)P(\texttt a\mid\texttt m)=\tfrac13$, lo mismo que la «a» inicial de «ama», $P(\texttt a\mid\fin)=\tfrac13$; el resto de ambas palabras es idéntico.
\end{enumerate}
\resultado \textbf{$P(\text{mar})=1/30\approx 0{,}0333$; $P(\text{rama})=2/75\approx 0{,}0267$; $P(\text{mama})=P(\text{ama})=4/75\approx 0{,}0533$.}
\comprobacion $P(\text{mama})/P(\text{ama})=\big(\tfrac13\cdot 1\big)/\tfrac13=1$.
\criterio 1 punto por probabilidad; se resta 1 punto si no se explica la generalización.

\apartado{(c)}
\datos $P(\text{mar})=1/30$, $P(\text{ama})=4/75$, $P(\text{rama})=2/75$; 13 predicciones; $\abs{\V}=4$.
\planteamiento pérdida media $\mathcal J=-\frac1{13}\sum\ln P(\text{palabra})$; en bits, $\mathcal J/\ln 2$; $\PPL=e^{\mathcal J}$ (M01, §5.1).
\pasos
\begin{enumerate}[label=\arabic*.]
  \item $\ln(1/30)=-3{,}4012$; $\ln(4/75)=-2{,}9312$; $\ln(2/75)=-3{,}6243$.
  \item Suma: $-9{,}9567$.
  \item $\mathcal J=9{,}9567/13=0{,}7659$ nats.
  \item En bits: $0{,}7659/0{,}6931=1{,}105$ bits.
  \item $\PPL=e^{0{,}7659}=2{,}151$.
  \item Modelo uniforme: $P=1/4$ en cada paso, $\mathcal J=\ln 4=1{,}386$ nats $=2$ bits y $\PPL=4$.
\end{enumerate}
\resultado \textbf{$\mathcal J\approx 0{,}766$ nats $\approx 1{,}10$ bits; $\PPL\approx 2{,}15$, frente a $4$ del modelo uniforme.} La comparación es optimista: se mide sobre el propio corpus de entrenamiento, no sobre datos de validación (M01, §5.3).
\comprobacion $\PPL=P(\text{corpus})^{-1/13}=\big(\tfrac1{30}\cdot\tfrac4{75}\cdot\tfrac2{75}\big)^{-1/13}=(4{,}741\cdot10^{-5})^{-1/13}=2{,}151$.
\criterio 1 punto la pérdida; 1 punto bits y perplejidad; 1 punto la comparación, con la advertencia sobre el conjunto de entrenamiento.

\apartado{(d)}
\datos fila \texttt m: $(0,3,0,0)$; fila \texttt r: $(1,1,0,0)$; $\abs{\V}=4$; $\alpha=1$.
\planteamiento $P_\alpha(v\mid u)=(N(u,v)+\alpha)/\big(\sum_w N(u,w)+\alpha\abs{\V}\big)$ (M01, §3.3).
\pasos
\begin{enumerate}[label=\arabic*.]
  \item Fila \texttt m suavizada: $(1,4,1,1)$, total $3+4=7$; $P(\texttt a\mid\texttt m)=4/7\approx 0{,}5714$.
  \item Fila \texttt r suavizada: $(2,2,1,1)$, total $2+4=6$; $P(\texttt r\mid\texttt r)=1/6\approx 0{,}1667$.
\end{enumerate}
\resultado \textbf{$P(\bin{a}\mid\bin{m})=4/7\approx 0{,}571$; $P(\bin{r}\mid\bin{r})=1/6\approx 0{,}167$.}
\comprobacion la fila \texttt m suavizada suma $(1+4+1+1)/7=1$; el suavizado reduce la probabilidad máxima ($1\to 0{,}571$) y da probabilidad positiva a los ceros.

\apartado{(e)}
\datos fila \texttt a: $p=(0{,}4;\ 0;\ 0{,}4;\ 0{,}2)$ en el orden $(\fin,\texttt a,\texttt m,\texttt r)$.
\planteamiento se calcula la acumulada $F_k$ y se elige el menor $k$ con $u<F_k$ (M01, §4.3).
\pasos
\begin{enumerate}[label=\arabic*.]
  \item $F=(0{,}4;\ 0{,}4;\ 0{,}8;\ 1)$.
  \item $u=0{,}35<F_1=0{,}4$: sale \fin.
  \item $u=0{,}55$: $F_1=F_2=0{,}4\le 0{,}55<F_3=0{,}8$: sale \texttt m. El token \texttt a, de probabilidad~0, ocupa un intervalo vacío y nunca sale.
  \item $u=0{,}9$: $0{,}8\le 0{,}9<1$: sale \texttt r.
\end{enumerate}
\resultado se obtienen, por este orden, los tokens \fin, \bin{m} y \bin{r}.
\comprobacion los intervalos $[0;\,0{,}4)$, $[0{,}4;\,0{,}8)$ y $[0{,}8;\,1)$ tienen longitudes $0{,}4$, $0{,}4$ y $0{,}2$, iguales a las probabilidades.

\repasar{M01, §2.3 y §3.2 (apartados a y b); M01, §5.1 y §5.3 (c); M01, §3.3 (d); M01, §4.3 (e).}
\end{solucion}

% =====================================================================  B.2
\begin{solucion}{eje:p2}
\apartado{(a)}
\datos $e_{\text{rey}}=(3,2,1)$, $e_{\text{reina}}=(3,-2,1)$.
\planteamiento producto escalar componente a componente, $\norm{v}=\sqrt{v\cdot v}$ y $\cos(u,v)=u\cdot v/(\norm u\norm v)$ (M02, §1.2, §1.3 y §2.3).
\pasos
\begin{enumerate}[label=\arabic*.]
  \item $e_{\text{rey}}\cdot e_{\text{reina}}=9-4+1=6$.
  \item $\norm{e_{\text{rey}}}=\sqrt{9+4+1}=\sqrt{14}\approx 3{,}742$, e igualmente $\norm{e_{\text{reina}}}=\sqrt{14}$.
  \item $\cos=6/14=3/7\approx 0{,}4286$.
\end{enumerate}
\resultado \textbf{producto escalar $6$; $\norm{e_{\text{rey}}}=\sqrt{14}\approx 3{,}74$; $\cos=3/7\approx 0{,}429$.}
\comprobacion $0<3/7<1$: los vectores coinciden en las coordenadas primera y tercera y difieren en el signo de la segunda.

\apartado{(b)}
\datos tabla de embeddings.
\planteamiento aritmética de vectores y coseno con cada candidato (M02, §5.4).
\pasos
\begin{enumerate}[label=\arabic*.]
  \item $q=(3-0+1,\ 2-2-2,\ 1-2+2)=(4,-2,1)$ y $\norm q=\sqrt{21}\approx 4{,}583$.
  \item Cosenos:
\end{enumerate}
\begin{center}
\small
\begin{tabular}{@{}lccc@{}}
\toprule
Palabra $w$ & $q\cdot e_w$ & $\norm{e_w}$ & $\cos(q,e_w)$ \\
\midrule
reina & $12+4+1=17$ & $\sqrt{14}$ & $17/\sqrt{294}\approx 0{,}991$ \\
princesa & $8+2+2=12$ & $3$ & $12/(3\sqrt{21})\approx 0{,}873$ \\
niña & $0+2+2=4$ & $\sqrt5$ & $4/\sqrt{105}\approx 0{,}390$ \\
\bottomrule
\end{tabular}
\end{center}
\begin{enumerate}[label=\arabic*.,start=3]
  \item Las palabras de la consulta tampoco ganan: rey, $9/\sqrt{294}\approx 0{,}525$; mujer, $10/(3\sqrt{21})\approx 0{,}727$; hombre, $-2/(\sqrt{21}\sqrt8)\approx -0{,}154$.
  \item La segunda coordenada vale $+2$ en rey y hombre y $-2$ en reina y mujer: codifica el género.
\end{enumerate}
\resultado \textbf{$q=(4,-2,1)$; la analogía devuelve «reina» ($\cos\approx 0{,}991$); el género está en la segunda coordenada.}
\comprobacion $q-e_{\text{reina}}=(1,0,0)$: $q$ solo difiere de reina en una unidad de la primera coordenada.
\criterio 1 punto $q$; 2 puntos los cosenos; 1 punto la interpretación.

\apartado{(c)}
\datos $W=\begin{pmatrix}1&0&0\\0&0&1\end{pmatrix}$.
\planteamiento $Wv=(v_1,v_3)$: $W$ conserva la primera y la tercera coordenadas (M02, §3.3); el núcleo se define en M03, §1.3.
\pasos
\begin{enumerate}[label=\arabic*.]
  \item $We_{\text{rey}}=(3,1)$ y $We_{\text{reina}}=(3,1)$.
  \item Coseno de las imágenes: $10/10=1$.
  \item $Wv=0\iff v_1=v_3=0$: el núcleo es la recta generada por $(0,1,0)$ y $\rango W=2$.
\end{enumerate}
\resultado \textbf{ambas imágenes son $(3,1)$, con coseno~1. El núcleo de $W$ es $\gen\{(0,1,0)\}$: $W$ destruye la coordenada del género.}
\comprobacion $e_{\text{rey}}-e_{\text{reina}}=(0,4,0)$ está en el núcleo, pues $W(0,4,0)\T=(0,0)\T$.

\apartado{(d)}
\datos vector de dimensión $1024$; matriz de $50\,000\times 1024$; $10^{12}$ operaciones por segundo.
\planteamiento cada logit es un producto escalar de $1024$ términos: $1024$ productos y $1023\approx 1024$ sumas (M02, §4.4).
\pasos
\begin{enumerate}[label=\arabic*.]
  \item Operaciones por logit: $2\cdot 1024=2048$.
  \item Por token: $50\,000\cdot 2048=1{,}024\cdot 10^{8}$.
  \item Tiempo: $1{,}024\cdot10^{8}/10^{12}=1{,}024\cdot 10^{-4}$ s $\approx 0{,}1$ ms.
\end{enumerate}
\resultado \textbf{$\approx 1{,}02\cdot 10^{8}$ operaciones por token, unos $0{,}102$ ms.}
\comprobacion regla de dos operaciones por peso: la matriz tiene $5{,}12\cdot 10^{7}$ pesos y $2\cdot 5{,}12\cdot10^{7}=1{,}024\cdot10^{8}$.
\criterio se acepta cualquier valor entre $5\cdot 10^{7}$ y $1{,}1\cdot 10^{8}$ con razonamiento correcto (contar solo productos da la mitad).

\repasar{M02, §1.2, §1.3 y §2.3 (a); M02, §5.4 (b); M02, §3.3 y M03, §1.3 (c); M02, §4.4 (d).}
\end{solucion}

% =====================================================================  B.3
\begin{solucion}{eje:p3}
\apartado{(a)}
\datos $A=\begin{pmatrix}3&1\\1&3\end{pmatrix}$.
\planteamiento polinomio característico $\det(A-\lambda I)=0$ y resolución de $(A-\lambda I)v=0$ (M03, §3.1).
\pasos
\begin{enumerate}[label=\arabic*.]
  \item $\det(A-\lambda I)=(3-\lambda)^{2}-1=\lambda^{2}-6\lambda+8=(\lambda-4)(\lambda-2)$.
  \item $\lambda_1=4$: $A-4I=\begin{pmatrix}-1&1\\1&-1\end{pmatrix}$ da $v_1=v_2$; vector unitario $q_1=(1,1)/\sqrt2$.
  \item $\lambda_2=2$: $A-2I=\begin{pmatrix}1&1\\1&1\end{pmatrix}$ da $v_1=-v_2$; vector unitario $q_2=(1,-1)/\sqrt2$.
  \item $q_1\cdot q_2=(1-1)/2=0$.
\end{enumerate}
\resultado \textbf{$\lambda=4$ con $(1,1)/\sqrt2$ y $\lambda=2$ con $(1,-1)/\sqrt2$. Son ortogonales porque $A$ es simétrica: por el teorema espectral (M03, §3.3), los autovectores de autovalores distintos de una matriz simétrica son ortogonales.}
\comprobacion $\tr A=6=4+2$ y $\det A=8=4\cdot 2$.
\criterio 1 punto el polinomio y los autovalores; 1 punto los autovectores unitarios; 1 punto la justificación.

\apartado{(b)}
\datos $\det A=8$.
\planteamiento para una matriz $2\times2$, $A^{-1}=\frac1{\det A}\begin{pmatrix}d&-b\\-c&a\end{pmatrix}$ (M03, §2.1).
\pasos
\begin{enumerate}[label=\arabic*.]
  \item $A^{-1}=\frac18\begin{pmatrix}3&-1\\-1&3\end{pmatrix}=\begin{pmatrix}0{,}375&-0{,}125\\-0{,}125&0{,}375\end{pmatrix}$.
  \item $AA^{-1}=\frac18\begin{pmatrix}9-1&-3+3\\3-3&-1+9\end{pmatrix}=I$.
  \item Si $Av=\lambda v$ con $\lambda\neq0$, entonces $A^{-1}v=\lambda^{-1}v$: $A^{-1}$ tiene los mismos autovectores y autovalores $1/4$ y $1/2$.
\end{enumerate}
\resultado \textbf{$A^{-1}=\frac18\begin{pmatrix}3&-1\\-1&3\end{pmatrix}$, con autovalores $1/4$ y $1/2$, inversos de los de $A$.}
\comprobacion $A^{-1}(1,1)\T=\tfrac18(2,2)\T=\tfrac14(1,1)\T$ y $A^{-1}(1,-1)\T=\tfrac18(4,-4)\T=\tfrac12(1,-1)\T$.

\apartado{(c)}
\datos $B=\begin{pmatrix}0&-2\\3&0\end{pmatrix}$.
\planteamiento los $v_i$ y $\sigma_i^{2}$ son los autovectores y autovalores de $B\T B$; las columnas de $U$ se obtienen como $u_i=Bv_i/\sigma_i$ (M03, §4.2). Por Eckart--Young, $B_1=\sigma_1u_1v_1\T$ y $\norm{B-B_1}_F=\sigma_2$ (M03, §5.1).
\pasos
\begin{enumerate}[label=\arabic*.]
  \item $B\T B=\begin{pmatrix}0&3\\-2&0\end{pmatrix}\begin{pmatrix}0&-2\\3&0\end{pmatrix}=\begin{pmatrix}9&0\\0&4\end{pmatrix}$.
  \item Es diagonal: $\sigma_1=\sqrt9=3$ con $v_1=(1,0)$, y $\sigma_2=\sqrt4=2$ con $v_2=(0,1)$; $V=I$ y $\Sigma=\diag(3,2)$.
  \item $u_1=Bv_1/3=(0,3)/3=(0,1)$ y $u_2=Bv_2/2=(-2,0)/2=(-1,0)$; $U=\begin{pmatrix}0&-1\\1&0\end{pmatrix}$, que es ortogonal.
  \item $B_1=3\,u_1v_1\T=3\begin{pmatrix}0&0\\1&0\end{pmatrix}=\begin{pmatrix}0&0\\3&0\end{pmatrix}$.
  \item $B-B_1=\begin{pmatrix}0&-2\\0&0\end{pmatrix}$ y $\norm{B-B_1}_F=2=\sigma_2$.
  \item Energía retenida: $9/(9+4)=9/13\approx 0{,}6923$.
  \item $B$ no es simétrica: su polinomio característico es $\lambda^{2}+6$, con raíces complejas $\pm i\sqrt6$. Sus autovalores no son reales ni sus autovectores ortogonales, de modo que no pueden servir de valores y vectores singulares. Solo para matrices simétricas semidefinidas positivas, como $A$, coinciden la SVD y la descomposición espectral.
\end{enumerate}
\resultado \textbf{$B=\begin{pmatrix}0&-1\\1&0\end{pmatrix}\begin{pmatrix}3&0\\0&2\end{pmatrix}\begin{pmatrix}1&0\\0&1\end{pmatrix}$; $B_1=\begin{pmatrix}0&0\\3&0\end{pmatrix}$; error $2$; energía retenida $9/13\approx 69{,}2\,\%$.}
\comprobacion $U\Sigma V\T=\begin{pmatrix}0\cdot3&-1\cdot2\\1\cdot3&0\cdot2\end{pmatrix}=B$, y $\norm{B}_F^{2}=4+9=13=\sigma_1^{2}+\sigma_2^{2}$.
\criterio 1 punto $V$ y $\Sigma$; 1 punto $U$ y la reconstrucción; 1 punto $B_1$, el error, la energía y la explicación.

\apartado{(d)}
\datos $m\times n=2048\times1024$; $k=64$; LoRA con $r=8$ sobre $d\times d=4096\times4096$.
\planteamiento almacenar el truncamiento de rango $k$ cuesta $k(m+n+1)$ números (M03, §5.2); LoRA entrena $B\in\R^{d_{\text{out}}\times r}$ y $A\in\R^{r\times d_{\text{in}}}$, es decir, $r(d_{\text{in}}+d_{\text{out}})$ parámetros (M03, §5.3).
\pasos
\begin{enumerate}[label=\arabic*.]
  \item Original: $2048\cdot1024=2\,097\,152$ números.
  \item Rango 64: $64\cdot(2048+1024+1)=64\cdot3073=196\,672$ números.
  \item Fracción: $196\,672/2\,097\,152=0{,}09378$, un $9{,}38\,\%$.
  \item Compensa mientras $3073\,k<2\,097\,152$, es decir, $k<682{,}4$: con $k=682$ se almacenan $2\,095\,786$ números y aún se ahorra; con $k=683$, $2\,098\,859$, que ya supera el original.
  \item LoRA: $8\cdot(4096+4096)=65\,536$ parámetros, frente a $4096^{2}=16\,777\,216$.
  \item Porcentaje: $65\,536/16\,777\,216=0{,}003906$, un $0{,}391\,\%$.
\end{enumerate}
\resultado \textbf{$196\,672$ números ($9{,}38\,\%$); compensa hasta $k=682$; LoRA de rango~8 entrena $65\,536$ parámetros, el $0{,}391\,\%$ de la matriz.}
\comprobacion $65\,536/16\,777\,216=2^{16}/2^{24}=2^{-8}=1/256\approx 0{,}39\,\%$.
\criterio 1 punto por cada cifra: almacenamiento, umbral, parámetros de LoRA y porcentaje.

\repasar{M03, §3.1 y §3.3 (a); M03, §2.1 y §3.2 (b); M03, §4.1, §4.2 y §5.1 (c); M03, §5.2 y §5.3 (d).}
\end{solucion}

% =====================================================================  B.4
\begin{solucion}{eje:p4}
\apartado{(a)}
\datos $L=-\ln p_c$, $p=\softmax(z)$, $S=\sum_j e^{z_j}$, $y$ one-hot de $c$.
\planteamiento reescribir la pérdida con la definición de la softmax y derivar término a término (M04, §3.4).
\pasos
\begin{enumerate}[label=\arabic*.]
  \item $H(y,p)=-\sum_k y_k\ln p_k=-\ln p_c$, porque $y_k=0$ para $k\neq c$ e $y_c=1$. Luego $L=H(y,p)$.
  \item $\ln p_c=\ln(e^{z_c}/S)=z_c-\ln S$, de donde $L=-z_c+\ln S$.
  \item $\partial(-z_c)/\partial z_k=-\mathbb I[k=c]=-y_k$.
  \item Regla de la cadena: $\partial(\ln S)/\partial z_k=\frac1S\,\partial S/\partial z_k=e^{z_k}/S=p_k$, pues solo el sumando $j=k$ de $S$ depende de $z_k$.
  \item Sumando: $\partial L/\partial z_k=p_k-y_k$. \qed
\end{enumerate}
\resultado \textbf{$L=H(y,p)=-z_c+\ln\sum_je^{z_j}$ y $\nabla_zL=p-y$.}
\comprobacion $\sum_k(p_k-y_k)=1-1=0$, coherente con la invariancia de la softmax por traslación: sumar una constante a todos los logits no cambia $L$.
\criterio 1 punto la equivalencia con la entropía cruzada; 1 punto la expresión con log-suma-exp; 2 puntos la derivación.

\apartado{(b)}
\datos $W$, $x=(1,2)$, $c=1$.
\planteamiento $z=Wx$, $p=\softmax(z)$, $L=-\ln p_1$.
\pasos
\begin{enumerate}[label=\arabic*.]
  \item $z=(1\cdot1+0\cdot2,\ 0\cdot1+1\cdot2,\ 1\cdot1-1\cdot2)=(1,2,-1)$.
  \item $e^{z}=(2{,}7183;\ 7{,}3891;\ 0{,}36788)$ y $S=10{,}475$.
  \item $p=(0{,}25950;\ 0{,}70538;\ 0{,}035119)$.
  \item $L=-\ln 0{,}25950=1{,}3490$ nats.
\end{enumerate}
\resultado \textbf{$z=(1,2,-1)$; $p\approx(0{,}260;\ 0{,}705;\ 0{,}0351)$; $L\approx 1{,}35$ nats.}
\comprobacion $p$ suma $1{,}0000$, y $L=-z_1+\ln S=-1+2{,}3490=1{,}3490$.

\apartado{(c)}
\datos $p$ del apartado (b), $y=(1,0,0)$, $x=(1,2)$.
\planteamiento $\delta=\nabla_zL=p-y$; para la capa lineal $z=Wx$, $\partial L/\partial W=\delta x\T$ y $\partial L/\partial x=W\T\delta$ (M04, §3.2).
\pasos
\begin{enumerate}[label=\arabic*.]
  \item $\delta=(-0{,}74050;\ 0{,}70538;\ 0{,}035119)$, de forma $3\times1$.
  \item $\partial L/\partial W=\delta x\T$, de forma $(3\times1)(1\times2)=3\times2$, como $W$:
  \[
  \frac{\partial L}{\partial W}=\begin{pmatrix}-0{,}7405&-1{,}4810\\0{,}7054&1{,}4108\\0{,}0351&0{,}0702\end{pmatrix}.
  \]
  \item $\partial L/\partial x=W\T\delta$, de forma $(2\times3)(3\times1)=2\times1$:
  \[
  W\T\delta=\begin{pmatrix}1&0&1\\0&1&-1\end{pmatrix}\begin{pmatrix}-0{,}7405\\0{,}7054\\0{,}0351\end{pmatrix}=\begin{pmatrix}-0{,}7054\\0{,}6703\end{pmatrix}.
  \]
\end{enumerate}
\resultado \textbf{$\nabla_zL\approx(-0{,}741;\ 0{,}705;\ 0{,}0351)$; $\partial L/\partial W$ como arriba ($3\times2$); $\partial L/\partial x\approx(-0{,}705;\ 0{,}670)$ ($2\times1$).}
\comprobacion las componentes de $\delta$ suman $0$, y también cada columna de $\partial L/\partial W$.
\criterio 1 punto $\delta$; 2 puntos $\partial L/\partial W$ con su forma; 1 punto $\partial L/\partial x$.

\apartado{(d)}
\datos $\eta=0{,}5$; $W$ y $\partial L/\partial W$ del apartado (c).
\planteamiento $W'=W-\eta\,\partial L/\partial W$ y nueva evaluación de $z'=W'x$, $p'$ y $L'$ (M04, §1.5).
\pasos
\begin{enumerate}[label=\arabic*.]
  \item $W'=\begin{pmatrix}1+0{,}37025&0+0{,}74050\\0-0{,}35269&1-0{,}70538\\1-0{,}01756&-1-0{,}03512\end{pmatrix}=\begin{pmatrix}1{,}3703&0{,}7405\\-0{,}3527&0{,}2946\\0{,}9824&-1{,}0351\end{pmatrix}$.
  \item $z'_1=1{,}3703+2\cdot0{,}7405=2{,}8513$; $z'_2=-0{,}3527+2\cdot0{,}2946=0{,}2365$; $z'_3=0{,}9824-2\cdot1{,}0351=-1{,}0878$.
  \item $e^{z'}=(17{,}310;\ 1{,}2669;\ 0{,}33696)$, con suma $18{,}913$.
  \item $p'_1=17{,}310/18{,}913=0{,}91520$ y $L'=-\ln0{,}91520=0{,}08861$ nats.
\end{enumerate}
\resultado \textbf{$W'\approx\begin{pmatrix}1{,}37&0{,}741\\-0{,}353&0{,}295\\0{,}982&-1{,}04\end{pmatrix}$; $p'_1\approx0{,}915$; la pérdida baja de $1{,}35$ a $0{,}0886$ nats.}
\comprobacion la aproximación de primer orden predice un descenso $\eta\norm{\partial L/\partial W}_F^{2}=0{,}5\cdot5{,}236=2{,}62$, mayor que el real, $1{,}260$: el paso es grande y la curvatura frena el descenso (M04, §4.4), pero la dirección es de descenso.

\apartado{(e)}
\datos $h=0{,}01$ sobre $W_{11}$; valor analítico $a=p_1-1=-0{,}7405035$.
\planteamiento $n=[L(W_{11}+h)-L(W_{11}-h)]/(2h)$ y error relativo $r(a,n)$ (M04, §6.1 y §6.3). Hacen falta 7 cifras porque se restan números muy próximos.
\pasos
\begin{enumerate}[label=\arabic*.]
  \item $W_{11}=1{,}01$: $z=(1{,}01;\ 2;\ -1)$ y $L_+=1{,}3416168$.
  \item $W_{11}=0{,}99$: $z=(0{,}99;\ 2;\ -1)$ y $L_-=1{,}3564268$.
  \item $n=(1{,}3416168-1{,}3564268)/0{,}02=-0{,}7405020$.
  \item $r=\abs{a-n}/\max(\abs a,\abs n)=1{,}54\cdot10^{-6}/0{,}7405=2{,}08\cdot10^{-6}$.
\end{enumerate}
\resultado \textbf{$n\approx-0{,}740502$ frente a $a\approx-0{,}740504$; $r\approx2{,}1\cdot10^{-6}$: el gradiente analítico es correcto.}
\comprobacion la diferencia es el error de truncamiento previsto, $\tfrac{h^{2}}{6}L'''$, con $L'''=p_1(1-p_1)(1-2p_1)=0{,}09243$ respecto de $W_{11}$: $10^{-4}\cdot0{,}09243/6=1{,}54\cdot10^{-6}$, que coincide con $n-a$. Con $h=0{,}01$ el truncamiento domina sobre el redondeo, por lo que no cabe esperar $r\le10^{-7}$.

\repasar{M04, §3.3 y §3.4 (a); M04, §3.2 (b y c); M04, §1.5 y §4.4 (d); M04, §6.1 y §6.3 (e).}
\end{solucion}

% =====================================================================  B.5
\begin{solucion}{eje:p5}
\apartado{(a)}
\datos $P(L=k)=(1-q)^{k-1}q$; observaciones $3,4,5,5,6,7$, con $n=6$ y $\sum_ik_i=30$.
\planteamiento log-verosimilitud, condición de primer orden y comprobación del signo de la segunda derivada (M05, §4.2).
\pasos
\begin{enumerate}[label=\arabic*.]
  \item $\ell(q)=\sum_i\big[(k_i-1)\ln(1-q)+\ln q\big]=\big(\sum_ik_i-n\big)\ln(1-q)+n\ln q$.
  \item $\ell'(q)=\dfrac nq-\dfrac{\sum_ik_i-n}{1-q}=0\iff n(1-q)=q\big(\sum_ik_i-n\big)\iff\hat q=\dfrac{n}{\sum_ik_i}$.
  \item $\ell''(q)=-\dfrac n{q^{2}}-\dfrac{\sum_ik_i-n}{(1-q)^{2}}<0$ en $(0,1)$: $\ell$ es estrictamente cóncava y $\hat q$ es su máximo global.
  \item $\hat q=6/30=0{,}2$.
  \item $P(L=3)=0{,}8^{2}\cdot0{,}2=0{,}128$ y $\E[L]=1/0{,}2=5$ letras.
\end{enumerate}
\resultado \textbf{$\hat q=n/\sum_ik_i=0{,}2$; $P(L=3)=0{,}128$; $\E[L]=5$ letras.}
\comprobacion la media muestral es $30/6=5=1/\hat q$: el EMV iguala la media teórica a la muestral.
\criterio 2 puntos la derivación (con la comprobación del máximo); 1 punto $\hat q$; 1 punto los dos valores.

\apartado{(b)}
\datos tres dados de 4 caras, uno de 6 y uno de 12; resultado observado: 5.
\planteamiento tabla de Bayes: a posteriori $\propto$ a priori $\times$ verosimilitud; la constante de normalización es la probabilidad predictiva a priori (M05, §1.3).
\pasos
\begin{enumerate}[label=\arabic*.]
  \item Tabla:
\end{enumerate}
\begin{center}
\small
\begin{tabular}{@{}lcccc@{}}
\toprule
Dado & A priori & $P(5\mid\text{dado})$ & Producto & A posteriori \\
\midrule
4 caras & $3/5$ & $0$ & $0$ & $0$ \\
6 caras & $1/5$ & $1/6$ & $1/30$ & $2/3$ \\
12 caras & $1/5$ & $1/12$ & $1/60$ & $1/3$ \\
\midrule
Total & $1$ & & $1/20$ & $1$ \\
\bottomrule
\end{tabular}
\end{center}
\begin{enumerate}[label=\arabic*.,start=2]
  \item Probabilidad predictiva a priori: $P(5)=\tfrac1{30}+\tfrac1{60}=\tfrac1{20}=0{,}05$.
  \item Dado el tipo de dado, los dos lanzamientos son independientes; por probabilidad total con las probabilidades a posteriori, $P(6\text{ en el segundo}\mid5)=\tfrac23\cdot\tfrac16+\tfrac13\cdot\tfrac1{12}=\tfrac4{36}+\tfrac1{36}=\tfrac5{36}\approx0{,}139$.
\end{enumerate}
\resultado \textbf{a posteriori $(0,\ 2/3,\ 1/3)$; $P(5)=1/20$; $P(6\text{ en el segundo lanzamiento}\mid5)=5/36\approx0{,}139$.}
\comprobacion las probabilidades a posteriori suman $2/3+1/3=1$. Antes de observar nada, $P(6)=\tfrac35\cdot0+\tfrac15\cdot\tfrac16+\tfrac15\cdot\tfrac1{12}=\tfrac1{20}=0{,}05$; tras el~5 sube a $5/36\approx0{,}139$, porque la observación descarta los dados de 4 caras, que no pueden dar un~6.
\criterio 1,5 puntos la tabla; 0,5 puntos la predictiva a priori; 1 punto la predictiva a posteriori.

\apartado{(c)}
\datos $\hat p=0{,}78$; $n=600$; $z_{0{,}975}=1{,}96$; semianchura objetivo $0{,}01$.
\planteamiento error típico $\sqrt{\hat p(1-\hat p)/n}$; intervalo $\hat p\pm1{,}96\cdot\mathrm{ET}$ por el teorema central del límite; para el tamaño, despejar $n$ de $1{,}96\sqrt{\hat p(1-\hat p)/n}=0{,}01$ (M05, §5.5).
\pasos
\begin{enumerate}[label=\arabic*.]
  \item $\hat p(1-\hat p)=0{,}78\cdot0{,}22=0{,}1716$.
  \item $\mathrm{ET}=\sqrt{0{,}1716/600}=\sqrt{2{,}86\cdot10^{-4}}=0{,}01691$.
  \item Semianchura: $1{,}96\cdot0{,}01691=0{,}03315$; intervalo $[0{,}7469;\ 0{,}8131]$.
  \item $n=(1{,}96/0{,}01)^{2}\cdot0{,}1716=38\,416\cdot0{,}1716=6592{,}2$, que se redondea hacia arriba: $6593$.
\end{enumerate}
\resultado \textbf{$\mathrm{ET}\approx0{,}0169$; intervalo del $95\,\%$ $\approx[74{,}7\,\%;\ 81{,}3\,\%]$; harían falta $6593$ preguntas.}
\comprobacion reducir la semianchura de $3{,}315$ a $1$ punto exige multiplicar $n$ por $3{,}315^{2}\approx10{,}99$, y $600\cdot10{,}99\approx6592$.
\criterio 1 punto el error típico; 1 punto el intervalo; 2 puntos el tamaño (se aceptan 6592 y 6593).

\repasar{M05, §2.2 y §4.2 (a); M05, §1.3 (b); M05, §5.4 y §5.5 (c).}
\end{solucion}

% =====================================================================  B.6
\begin{solucion}{eje:p6}
\apartado{(a)}
\datos $p=(1/4,1/4,1/4,1/8,1/8)$ para \texttt a, \texttt b, \texttt c, \texttt d, \texttt e; mensaje «\texttt{abed}».
\planteamiento $H(p)=-\sum_ip_i\log_2p_i$ (M06, §1.2); algoritmo de Huffman, que fusiona en cada paso los dos nodos de menor probabilidad (M06, §3.4); desigualdad de Kraft $\sum_i2^{-\ell_i}\le1$ (M06, §3.2).
\pasos
\begin{enumerate}[label=\arabic*.]
  \item $H(p)=3\cdot\tfrac14\cdot2+2\cdot\tfrac18\cdot3=1{,}5+0{,}75=2{,}25$ bits.
  \item Huffman: se fusionan \texttt d y \texttt e ($1/4$); quedan cuatro nodos de $1/4$, que se fusionan por pares, $\{\texttt a,\texttt b\}$ y $\{\texttt c,\texttt{de}\}$ ($1/2$ cada uno), y estos en la raíz.
  \item Código: $\texttt a=\bin{00}$, $\texttt b=\bin{01}$, $\texttt c=\bin{10}$, $\texttt d=\bin{110}$, $\texttt e=\bin{111}$. Cualquier otra elección en los empates da las mismas longitudes $(2,2,2,3,3)$.
  \item Longitud media: $3\cdot\tfrac14\cdot2+2\cdot\tfrac18\cdot3=2{,}25$ bits $=H(p)$.
  \item Kraft: $3\cdot2^{-2}+2\cdot2^{-3}=0{,}75+0{,}25=1$.
  \item «\texttt{abed}» $=\bin{00}\,|\,\bin{01}\,|\,\bin{111}\,|\,\bin{110}=\bin{0001111110}$: $2+2+3+3=10$ bits.
\end{enumerate}
\resultado \textbf{$H(p)=2{,}25$ bits; código $\{\bin{00},\bin{01},\bin{10},\bin{110},\bin{111}\}$ con longitud media $2{,}25$ bits $=H(p)$ y suma de Kraft~1; «\bin{abed}» $=\bin{0001111110}$, 10 bits.}
\comprobacion la longitud media alcanza la entropía porque todas las probabilidades son potencias de $1/2$ y $\ell_i=-\log_2p_i$ (M06, §3.3); además, $-\log_2P(\text{«abed»})=2+2+3+3=10$ bits.
\criterio 1 punto $H(p)$; 1 punto el código y la comprobación de Kraft; 1 punto el mensaje.

\apartado{(b)}
\datos $p$ como en (a); $q=(1/5,\dots,1/5)$; $q'=(1/8,1/8,1/4,1/4,1/4)$.
\planteamiento $H(p,q)=-\sum_ip_i\log_2q_i$ y $\KL{p}{q}=H(p,q)-H(p)$ (M06, §2.2 y §2.3).
\pasos
\begin{enumerate}[label=\arabic*.]
  \item $H(p,q)=\log_25=2{,}3219$ bits y $\KL{p}{q}=2{,}3219-2{,}25=0{,}0719$ bits.
  \item $H(p,q')=\tfrac14\cdot3+\tfrac14\cdot3+\tfrac14\cdot2+\tfrac18\cdot2+\tfrac18\cdot2=2{,}5$ bits y $\KL{p}{q'}=2{,}5-2{,}25=0{,}25$ bits.
\end{enumerate}
\resultado \textbf{$H(p,q)\approx2{,}32$ bits y $\KL{p}{q}\approx0{,}0719$ bits; $H(p,q')=2{,}5$ bits y $\KL{p}{q'}=0{,}25$ bits. Es mejor el modelo uniforme: $q'$ concentra la probabilidad en símbolos equivocados, y un modelo seguro y equivocado puede ser peor que uno que no sabe nada.}
\comprobacion cálculo directo: $\KL{p}{q'}=\tfrac14\log_22+\tfrac14\log_22+\tfrac14\log_21+\tfrac18\log_2\tfrac12+\tfrac18\log_2\tfrac12=0{,}25+0{,}25+0-0{,}125-0{,}125=0{,}25$.
\criterio 0,75 puntos por cada valor.

\apartado{(c)}
\datos $\mathcal J=1{,}8$ nats por token; $\bar b=4{,}2$ bytes por token; $10^{9}$ bytes.
\planteamiento $\PPL=e^{\mathcal J}$; bits por token $=\mathcal J/\ln2$; $\bpb=\mathcal J/(\bar b\ln2)$; tamaño comprimido $=B\cdot\bpb/8$ bytes (M06, §4.3 y §4.4).
\pasos
\begin{enumerate}[label=\arabic*.]
  \item $\PPL=e^{1{,}8}=6{,}050$.
  \item Bits por token: $1{,}8/0{,}69315=2{,}5969$.
  \item Bits por byte: $2{,}5969/4{,}2=0{,}61830$.
  \item Tamaño comprimido: $10^{9}\cdot0{,}61830/8=7{,}729\cdot10^{7}$ bytes $\approx77{,}3$ MB.
  \item Razón de compresión: $8/0{,}61830=12{,}94$.
  \item La codificación aritmética emplea menos de $-\log_2q(x_{1:N})+2$ bits para todo el texto (M06, §4.2): el exceso es despreciable.
\end{enumerate}
\resultado \textbf{$\PPL\approx6{,}05$; $2{,}60$ bits por token; $0{,}618$ bits por byte; 1~GB se comprime a unos $77{,}3$ MB, con una razón de compresión $\approx12{,}9$.}
\comprobacion $10^{9}/12{,}94=7{,}73\cdot10^{7}$ bytes, como en el paso~4.
\criterio 1 punto por cada cifra; se aceptan redondeos razonables.

\repasar{M06, §1.2, §3.2, §3.3 y §3.4 (a); M06, §2.2 y §2.3 (b); M06, §4.2, §4.3 y §4.4, y M01, §5.1 (c).}
\end{solucion}
